% ---------------------------------------------------------------------------
% Study report · The thermomechanical response of the wall
%
% Build:  pdflatex thermomechanical_response_report.tex  (twice, for the ToC)
% Figures and table bodies are read from ../outputs/ and are produced by the
% paired notebook, which must be run first.
% ---------------------------------------------------------------------------
\documentclass[11pt,a4paper]{article}

\newcommand{\runninghead}{The thermomechanical response at Gubbio}

% ---------------------------------------------------------------------------
% reportstyle.tex — shared preamble for the study reports under studies/.
%
% Kept deliberately inside the package set shipped with TeX Live "basic":
% no tcolorbox, no siunitx, no enumitem, no titlesec. Everything below
% compiles with a stock pdflatex installation.
%
% Usage, from a report directory one level below a study folder:
%     % ---------------------------------------------------------------------------
% reportstyle.tex — shared preamble for the study reports under studies/.
%
% Kept deliberately inside the package set shipped with TeX Live "basic":
% no tcolorbox, no siunitx, no enumitem, no titlesec. Everything below
% compiles with a stock pdflatex installation.
%
% Usage, from a report directory one level below a study folder:
%     % ---------------------------------------------------------------------------
% reportstyle.tex — shared preamble for the study reports under studies/.
%
% Kept deliberately inside the package set shipped with TeX Live "basic":
% no tcolorbox, no siunitx, no enumitem, no titlesec. Everything below
% compiles with a stock pdflatex installation.
%
% Usage, from a report directory one level below a study folder:
%     \input{../../_shared/reportstyle.tex}
%     \graphicspath{{../outputs/}}
% ---------------------------------------------------------------------------

\usepackage[T1]{fontenc}
\usepackage[utf8]{inputenc}
\usepackage{textcomp}
\usepackage{lmodern}
\usepackage{microtype}
\usepackage[margin=2.4cm,headheight=15pt]{geometry}
\usepackage{graphicx}
\usepackage{booktabs}
\usepackage{tabularx}
\usepackage{array}
\usepackage{amsmath}
\usepackage{xcolor}
\usepackage[font=small,labelfont=bf,justification=justified]{caption}
\usepackage{float}
\usepackage{fancyhdr}
\usepackage[hidelinks]{hyperref}

% --- Palette ---------------------------------------------------------------
\definecolor{rulegrey}{HTML}{B9C0C7}
\definecolor{fillgrey}{HTML}{F2F4F6}
\definecolor{failred}{HTML}{A5202B}
\definecolor{passgreen}{HTML}{1B6B3A}
\definecolor{accent}{HTML}{1F4E79}

% --- Semantic shorthands ---------------------------------------------------
\newcommand{\degC}{\textdegree C}
\newcommand{\mdegC}{mdeg/\textdegree C}
\newcommand{\fail}{\textcolor{failred}{\textbf{fail}}}
\newcommand{\pass}{\textcolor{passgreen}{\textbf{pass}}}
\newcommand{\worse}{\textcolor{failred}{worse}}
\newcommand{\better}{\textcolor{passgreen}{better}}

% --- Highlighted panel -----------------------------------------------------
% #1 = heading, #2 = body
\newcommand{\panel}[2]{%
  \begin{center}%
  \setlength{\fboxsep}{9pt}%
  \fcolorbox{rulegrey}{fillgrey}{%
    \begin{minipage}{0.93\textwidth}%
      {\bfseries\color{accent}#1}\par\medskip
      #2
    \end{minipage}}%
  \end{center}}

% --- Numeric column: right aligned, fixed width ----------------------------
\newcolumntype{R}[1]{>{\raggedleft\arraybackslash}p{#1}}
\newcolumntype{L}[1]{>{\raggedright\arraybackslash}p{#1}}
\newcolumntype{Y}{>{\raggedright\arraybackslash}X}

% --- Table look ------------------------------------------------------------
\renewcommand{\arraystretch}{1.15}
\setlength{\tabcolsep}{6pt}

% --- Headers ---------------------------------------------------------------
\pagestyle{fancy}
\fancyhf{}
\fancyhead[L]{\small\itshape\runninghead}
\fancyhead[R]{\small\thepage}
\renewcommand{\headrulewidth}{0.4pt}

% --- Section spacing without titlesec --------------------------------------
\makeatletter
\renewcommand\section{\@startsection{section}{1}{\z@}%
  {-3.2ex \@plus -1ex \@minus -.2ex}{2.0ex \@plus.2ex}%
  {\normalfont\Large\bfseries\color{accent}}}
\renewcommand\subsection{\@startsection{subsection}{2}{\z@}%
  {-2.6ex \@plus -1ex \@minus -.2ex}{1.4ex \@plus.2ex}%
  {\normalfont\large\bfseries}}
\makeatother

% --- Figure defaults -------------------------------------------------------
\setkeys{Gin}{width=\linewidth}
\captionsetup[figure]{skip=6pt}

%     \graphicspath{{../outputs/}}
% ---------------------------------------------------------------------------

\usepackage[T1]{fontenc}
\usepackage[utf8]{inputenc}
\usepackage{textcomp}
\usepackage{lmodern}
\usepackage{microtype}
\usepackage[margin=2.4cm,headheight=15pt]{geometry}
\usepackage{graphicx}
\usepackage{booktabs}
\usepackage{tabularx}
\usepackage{array}
\usepackage{amsmath}
\usepackage{xcolor}
\usepackage[font=small,labelfont=bf,justification=justified]{caption}
\usepackage{float}
\usepackage{fancyhdr}
\usepackage[hidelinks]{hyperref}

% --- Palette ---------------------------------------------------------------
\definecolor{rulegrey}{HTML}{B9C0C7}
\definecolor{fillgrey}{HTML}{F2F4F6}
\definecolor{failred}{HTML}{A5202B}
\definecolor{passgreen}{HTML}{1B6B3A}
\definecolor{accent}{HTML}{1F4E79}

% --- Semantic shorthands ---------------------------------------------------
\newcommand{\degC}{\textdegree C}
\newcommand{\mdegC}{mdeg/\textdegree C}
\newcommand{\fail}{\textcolor{failred}{\textbf{fail}}}
\newcommand{\pass}{\textcolor{passgreen}{\textbf{pass}}}
\newcommand{\worse}{\textcolor{failred}{worse}}
\newcommand{\better}{\textcolor{passgreen}{better}}

% --- Highlighted panel -----------------------------------------------------
% #1 = heading, #2 = body
\newcommand{\panel}[2]{%
  \begin{center}%
  \setlength{\fboxsep}{9pt}%
  \fcolorbox{rulegrey}{fillgrey}{%
    \begin{minipage}{0.93\textwidth}%
      {\bfseries\color{accent}#1}\par\medskip
      #2
    \end{minipage}}%
  \end{center}}

% --- Numeric column: right aligned, fixed width ----------------------------
\newcolumntype{R}[1]{>{\raggedleft\arraybackslash}p{#1}}
\newcolumntype{L}[1]{>{\raggedright\arraybackslash}p{#1}}
\newcolumntype{Y}{>{\raggedright\arraybackslash}X}

% --- Table look ------------------------------------------------------------
\renewcommand{\arraystretch}{1.15}
\setlength{\tabcolsep}{6pt}

% --- Headers ---------------------------------------------------------------
\pagestyle{fancy}
\fancyhf{}
\fancyhead[L]{\small\itshape\runninghead}
\fancyhead[R]{\small\thepage}
\renewcommand{\headrulewidth}{0.4pt}

% --- Section spacing without titlesec --------------------------------------
\makeatletter
\renewcommand\section{\@startsection{section}{1}{\z@}%
  {-3.2ex \@plus -1ex \@minus -.2ex}{2.0ex \@plus.2ex}%
  {\normalfont\Large\bfseries\color{accent}}}
\renewcommand\subsection{\@startsection{subsection}{2}{\z@}%
  {-2.6ex \@plus -1ex \@minus -.2ex}{1.4ex \@plus.2ex}%
  {\normalfont\large\bfseries}}
\makeatother

% --- Figure defaults -------------------------------------------------------
\setkeys{Gin}{width=\linewidth}
\captionsetup[figure]{skip=6pt}

%     \graphicspath{{../outputs/}}
% ---------------------------------------------------------------------------

\usepackage[T1]{fontenc}
\usepackage[utf8]{inputenc}
\usepackage{textcomp}
\usepackage{lmodern}
\usepackage{microtype}
\usepackage[margin=2.4cm,headheight=15pt]{geometry}
\usepackage{graphicx}
\usepackage{booktabs}
\usepackage{tabularx}
\usepackage{array}
\usepackage{amsmath}
\usepackage{xcolor}
\usepackage[font=small,labelfont=bf,justification=justified]{caption}
\usepackage{float}
\usepackage{fancyhdr}
\usepackage[hidelinks]{hyperref}

% --- Palette ---------------------------------------------------------------
\definecolor{rulegrey}{HTML}{B9C0C7}
\definecolor{fillgrey}{HTML}{F2F4F6}
\definecolor{failred}{HTML}{A5202B}
\definecolor{passgreen}{HTML}{1B6B3A}
\definecolor{accent}{HTML}{1F4E79}

% --- Semantic shorthands ---------------------------------------------------
\newcommand{\degC}{\textdegree C}
\newcommand{\mdegC}{mdeg/\textdegree C}
\newcommand{\fail}{\textcolor{failred}{\textbf{fail}}}
\newcommand{\pass}{\textcolor{passgreen}{\textbf{pass}}}
\newcommand{\worse}{\textcolor{failred}{worse}}
\newcommand{\better}{\textcolor{passgreen}{better}}

% --- Highlighted panel -----------------------------------------------------
% #1 = heading, #2 = body
\newcommand{\panel}[2]{%
  \begin{center}%
  \setlength{\fboxsep}{9pt}%
  \fcolorbox{rulegrey}{fillgrey}{%
    \begin{minipage}{0.93\textwidth}%
      {\bfseries\color{accent}#1}\par\medskip
      #2
    \end{minipage}}%
  \end{center}}

% --- Numeric column: right aligned, fixed width ----------------------------
\newcolumntype{R}[1]{>{\raggedleft\arraybackslash}p{#1}}
\newcolumntype{L}[1]{>{\raggedright\arraybackslash}p{#1}}
\newcolumntype{Y}{>{\raggedright\arraybackslash}X}

% --- Table look ------------------------------------------------------------
\renewcommand{\arraystretch}{1.15}
\setlength{\tabcolsep}{6pt}

% --- Headers ---------------------------------------------------------------
\pagestyle{fancy}
\fancyhf{}
\fancyhead[L]{\small\itshape\runninghead}
\fancyhead[R]{\small\thepage}
\renewcommand{\headrulewidth}{0.4pt}

% --- Section spacing without titlesec --------------------------------------
\makeatletter
\renewcommand\section{\@startsection{section}{1}{\z@}%
  {-3.2ex \@plus -1ex \@minus -.2ex}{2.0ex \@plus.2ex}%
  {\normalfont\Large\bfseries\color{accent}}}
\renewcommand\subsection{\@startsection{subsection}{2}{\z@}%
  {-2.6ex \@plus -1ex \@minus -.2ex}{1.4ex \@plus.2ex}%
  {\normalfont\large\bfseries}}
\makeatother

% --- Figure defaults -------------------------------------------------------
\setkeys{Gin}{width=\linewidth}
\captionsetup[figure]{skip=6pt}

\graphicspath{{../outputs/}}

\title{\vspace{-1.2cm}\textbf{The thermomechanical response: what moves the wall,
and how long it takes}\\[2mm] \large A lag-and-gain screen of every candidate
driver, against the on-structure sensors, ERA5 and the Gubbio station}
\author{Study report · \texttt{studies/03\_thermomechanical\_response/}}
\date{}

\begin{document}
\maketitle
\thispagestyle{fancy}

\begin{abstract}
\noindent
Study~01 established what the monitoring record contains and Study~02
characterised the external sources that describe the weather at the site.
Neither measured how the wall itself responds. This report screens every
candidate driver the project has against one response---the compensated,
cleaned inclination---in three rounds, one for each source family: the sensors
on the structure, the ERA5 reanalysis, and the meteorological station in Gubbio
town. Each pair is scanned over the operators its physics admits---a transport
delay and a thermal time constant, jointly, because a filter and a shift can
produce the same peak correlation while meaning different things---and fitted at
the operator the scan selects, on the levels and on the diurnal band separately,
with Newey--West standard errors and a negative control carried through every
table.
Air temperature is the strongest driver in every round and its association is
negative, as the site geometry predicts; the response follows the on-structure
air temperature with no resolvable delay at hourly resolution, follows solar
radiation by one hour, and \emph{leads} the wall-temperature probe by two. The
screen is reported with its own noise floor, which turns out not to be silent,
and with the time constants it could not identify, which turn out to be most of
them.
\end{abstract}

\tableofcontents

% ===========================================================================
\section{Introduction}
\label{sec:introduction}

A grey-box model of this structure cannot be specified until three things are
known about each candidate driver: whether the wall responds to it at all, after
how long, and in which direction. This report measures those three quantities.

The question is not new to the project, but the answers on record are not
usable. Section~7.5 of \texttt{docs/raw-data-format.md} states the expectation
that the site geometry imposes---the valley face of the embankment is the more
exposed, so daytime heating expands it and tips the wall towards the mountain,
which is a \emph{negative} change by the instrument's convention---and records
against it a measured positive slope on the raw channel, on two different units
of hardware. Both statements come from studies now retired to
\texttt{studies/obsolete/}, whose conclusions this project does not build upon.

This study therefore re-measures the coupling on trusted ground: the response
and the on-structure drivers come from Study~01's exported archive, the external
drivers are read through the same loaders and channel maps Study~02 uses, and
every function that computes anything lives in \texttt{studies/shmlib/}.

\paragraph{What this report does not do.} It does not re-open the raw-versus-%
compensated contradiction of Section~7.5. The response screened here is
\texttt{inc\_comp\_cleaned}, Study~01's analysis column, and nothing else. The
contradiction remains recorded where it is recorded; the numbers below are
measured on the compensated series and describe that series.

% ===========================================================================
\section{Method}
\label{sec:method}

\subsection{The response}

\texttt{inc\_comp\_cleaned} on the hourly UTC grid. Three of Study~01's
decisions are consumed and none is revisited: the compensation, the join across
the February~2025 instrument change, and the spike verdict, which marks the
samples Study~01 replaced by interpolation. Those samples are dropped before any
hour is averaged. An interpolated value is the interpolator's output rather than
the instrument's, and a coupling measured against one would be partly a coupling
to Study~01's cleaning.

The absolute level of this column means nothing---it is fixed by how the
instrument sat in its mount and by the anchors Study~01 applied---so only its
changes carry structural information.

\subsection{The two bands}

Every series is screened twice: as its level, and as its diurnal band, defined
as the series minus its 24-hour centred rolling mean. The two answer different
questions. A correlation on the levels is dominated by whatever slow variation
the two series share; a correlation on the band is a statement about the daily
forcing cycle alone. The filter is one line of arithmetic rather than a designed
band-pass, so that its definition is reproducible and no unstated phase response
sits between the data and the lags measured afterwards. Figure~\ref{fig:bands}
shows what the separation does to one week of the response.

\begin{figure}[H]
\centering
\includegraphics[width=\textwidth]{TR_F01_response_bands.png}
\caption{The response over one week, as recorded and after the band separation.
The upper panel carries the level, whose position on the axis is set by an
arbitrary anchor; the lower panel carries the diurnal band, in which only the
daily excursion survives.}
\label{fig:bands}
\end{figure}

\subsection{Two things can delay a response, and they are not the same thing}

A \emph{transport delay} shifts a driver in time without changing its shape: it
is the time a thermal front takes to travel from the exposed face to the depth
that governs the movement. A \emph{thermal inertia} low-passes it: masonry does
not follow the air around it, it integrates it, and the discrete form of that
lumped-capacitance response is a one-pole filter with a time constant $\tau$.

The two are not interchangeable, and a scan over delay alone cannot tell them
apart: a filter with the right time constant and no delay, and a delay with no
filter, produce very similar peak correlations while implying entirely different
physics. Every pair in this report is therefore scanned over both, and the whole
grid is kept, because the maximum of such a grid is usually not a peak but a
ridge---delay and time constant substitute for one another over a wide range,
and a single winning cell quoted alone claims a resolution the record does not
have.

Two restrictions on that grid are load-bearing.

\paragraph{The diurnal band is scanned over delay alone.} On a band holding
essentially one period the two parameters are not separately identifiable at
all. A one-pole filter applied to a signal of one frequency is exactly an
amplitude scaling and a phase shift, and a delay is also a phase shift, so every
combination on a line of constant total phase fits equally well. Scanned jointly
there, the search rotates the phase past a half cycle and reports the flipped
sign as its optimum: the two external air temperatures, which stand at
$r = -0.81$ and $-0.79$ against the response, came back as $+0.82$ at a delay of
five hours and $\tau$ of eighteen. That is an artefact of fitting an
unidentifiable parameter, not a finding, and pinning the band's grid to
$\tau = 0$ removes it.

\paragraph{On the levels the grid stops at three days, by a stated prior.} The
mass whose temperature governs the inclination is a wall face and its effective
time constant is hours to a couple of days. Left unbounded the scan does not
settle on anything of that order: run out to 2160~h, the optimum for the
external drivers slides to between 720 and 1440~h---a month to two
months---at which point the filtered driver is no longer weather but season, and
what it correlates with is the response's own annual drift. The negative control
slides the same way. Section~\ref{sec:operator} reports which time constants
survived that bound and which did not.

\subsection{Lags, and the bound on them}

Each pair is scanned over integer-hour delays, the driver shifted forward so
that a positive lag means the response follows. The scan maximises $|r|$ rather than
$r$: the expected association here is negative, and a search maximising the
signed correlation would walk away from the optimum it was sent to find.

External forcings are scanned over $0$ to $12$~hours. Both bounds are physical.
A forcing must precede the response it causes, so a negative optimum against one
is not a measurement of anything. Twelve hours is where a delay stops being
distinguishable from a lead: on a near-periodic daily signal a delay of $d$
hours and a lead of $24-d$ produce the same alignment, and a lead cannot be
represented by a causal filter or used by a forecast.

Wall temperature is exempt and is scanned over $\pm 24$~hours. It is not a
forcing but an internal state variable at an unknown depth, and nothing requires
it to precede the deformation: a shallow probe would be led by the wall's motion
and a deep one would follow it.

\subsection{Gains, and their uncertainty}

At the selected operator---both the delay and the time constant---the response
is regressed on the driver by ordinary least squares, and the slope is reported
in millidegrees per unit of the driver. Fitting against the operator that won,
rather than against the driver as recorded, is not a detail: on a synthetic pair
built with a known gain of $-2.5$ and a six-hour time constant, regressing on
the unfiltered driver returns $-0.78$. Its
standard error is Newey--West with a one-day truncation. This is not a
refinement: the residuals of an hourly geophysical regression are strongly
autocorrelated, and the ordinary standard error of such a fit understates the
uncertainty by enough to make a coincidence look significant.

\subsection{The strata, and the negative control}

Every scan is run over the whole record, over each instrument era, and over each
season, on the season definitions Studies~01 and~02 use. Supply voltage is
carried through every round as a negative control. A bounded search over a
near-periodic pair returns an optimum for every candidate, so what separates a
coupling from that background is not the size of $|r|$ alone but its size
relative to a channel that has no mechanism by which it could move a wall.
Section~\ref{sec:control} reports what the control actually did, which is not
what was expected of it.

% ===========================================================================
\section{What the record supports}
\label{sec:coverage}

The three source families do not cover the same record. The on-structure air
temperature, humidity and supply voltage run the length of the post-outage archive; solar
radiation and wall temperature exist only in the current instrument era and are
further reduced by Study~01's verdicts, leaving roughly five thousand and
roughly forty-eight hundred paired hours respectively against some
twenty-one thousand for air temperature. Every result below is reported with
the paired hours it rests on, and the two current-era channels are read in that
light.

\begin{table}[H]
\centering
\small
\begin{tabular}{lrrll}
\toprule
\textbf{Driver} & \textbf{Hours} & \textbf{Paired} & \textbf{First} & \textbf{Last} \\
\midrule
\input{../outputs/TR_T01_driver_coverage.tex}\\
\bottomrule
\end{tabular}
\caption{What each driver has to work with: the hours it carries, the hours it
shares with the response, and the extent of that overlap.}
\label{tab:coverage}
\end{table}

% ===========================================================================
\section{Round 1 · the on-structure drivers}
\label{sec:round1}

The channels measured at the wall are the closest thing to a local forcing the
project has, and each carries a caveat Study~01 established: the air temperature
is recorded inside a sun-exposed housing rather than a meteorological screen,
the radiation exists only on the current-era days Study~01 did not condemn, and
the wall-temperature probe is absent on more than half the days of that era.

On the diurnal band the housing air temperature is the strongest association in
the whole study, $r = -0.957$ at zero lag, with a gain of
$-2.79$~\mdegC. Solar radiation follows at $r = -0.867$ with a delay of one hour
and a gain of $-0.026$~mdeg per watt per square metre. Both signs are the sign the
site's geometry predicts, and neither needs an operator: their time constants
are zero and the instantaneous fit is what the winning cell reproduces.

On the levels the on-structure radiation behaves differently. Its instantaneous
fit explains almost nothing, $R^2 = 0.17$, and a 48-hour time constant lifts that
to $0.29$ --- $\Delta R^2 = 0.122$, the largest operator gain of any driver whose
optimum is not against a bound. Read with Section~\ref{sec:operator}, that is a
statement about the slow band being dominated by variation neither the wall nor
the radiation resolves, rather than a measurement of the wall's thermal mass.

The wall temperature is the interesting case. On the diurnal band its optimum
sits at $-2$~hours and on the levels at $-4$~hours with a time constant of four,
and both say the same thing: the inclination \emph{leads} the probe. Under the
reading of Section~7.5 that is the signature of a probe sitting deep enough that
the layers actually driving the movement---those nearest the sun-exposed valley
face---reach their extremes before it does. It is a measurement, not an error.

It is also the one driver whose operator is worth what it costs. On the band the
delay buys $\Delta R^2 = 0.214$ over the instantaneous fit; on the levels the
delay and the time constant together buy $0.071$, and both parameters land
inside their grids rather than against a bound, so this is the only external
driver in the study whose thermal time constant the record actually identifies.

What the clamp costs can now be stated rather than asserted. Restricting the
search to non-negative delays---the restriction any forecast is under, since a
lead would need the probe's future values at the moment of issue---gives up
$\Delta R^2 = 0.068$ on the band and $0.024$ on the levels. That is the price of
using this channel predictively, and it is small enough to pay.

\begin{figure}[H]
\centering
\includegraphics[width=\textwidth]{TR_F02_lag_curves_str.png}
\caption{Round~1. Correlation against lag for each on-structure driver, diurnal
band, current era. The marked point on each curve is the optimum the scan
selected. The supply voltage is drawn with the rest: it is the negative control,
and its curve is the floor the others are read against.}
\label{fig:lags_str}
\end{figure}

% ===========================================================================
\section{Round 2 · ERA5}
\label{sec:round2}

ERA5 is a nine-kilometre grid average: its air temperature is not the air
temperature at the wall and its radiation is not the radiation the wall
receives. What it offers instead is a record with no gaps across the whole
archive, which is precisely what the on-structure channels lack.

On the diurnal band its air temperature reaches $r = -0.787$ at zero lag with a
gain of $-2.04$~\mdegC, and its radiation $r = -0.873$, also at zero lag. Both
are weaker than their on-structure counterparts and both keep the expected sign.
The attenuation is what a spatial average should do to a local signal, and it is
consistent with Study~02's finding that the three sources are not replicates of
one quantity.

\begin{figure}[H]
\centering
\includegraphics[width=\textwidth]{TR_F03_lag_curves_era5.png}
\caption{Round~2. Correlation against lag for the ERA5 channels, diurnal band,
whole record.}
\label{fig:lags_era5}
\end{figure}

% ===========================================================================
\section{Round 3 · the town station}
\label{sec:round3}

The Gubbio station is a point observation from a standard screen, neither a grid
average nor a box bolted to the wall. On the diurnal band its radiation gives
the strongest external association of the study, $r = -0.879$ at zero lag with a
gain of $-0.035$~mdeg per watt per square metre, and its air temperature
$r = -0.807$ with a gain of $-2.23$~\mdegC. Both sit between the on-structure
channels and ERA5, which is where a town-centre screen ought to sit.

\begin{figure}[H]
\centering
\includegraphics[width=\textwidth]{TR_F04_lag_curves_gs.png}
\caption{Round~3. Correlation against lag for the Gubbio station channels,
diurnal band, whole record.}
\label{fig:lags_gs}
\end{figure}

% ===========================================================================
\section{What the operator scan settled, and what it did not}
\label{sec:operator}

Scanning delay and time constant jointly was meant to answer whether the wall's
response is delayed, damped, or neither. It answered that for two drivers and
refused to answer it for the rest, and the refusal is the more useful half.

\paragraph{Two time constants are identified.} The wall temperature settles at
$\tau = 4$~h on the levels and the on-structure radiation at $\tau = 48$~h, both
strictly inside the grid, both buying real explanatory power over the
instantaneous fit ($\Delta R^2 = 0.071$ and $0.122$). The radiation's placement
is the weaker of the two: its surface is nearly flat from 36 to 72~h
($R^2 = 0.285$, $0.288$, $0.287$), so what that panel establishes is that the
coupling is inertial rather than instantaneous, not that the time constant is
two days. On the diurnal band every
driver's time constant is zero by construction, and the delays are zero or one
hour, so the daily coupling is instantaneous at hourly resolution.

\paragraph{Every other time constant on the levels is unidentified.} Fifteen
of the twenty-two drivers pin against the three-day bound rather than settling
inside it, which means the scan stopped where the grid did and not where the
physics did. Figures~\ref{fig:operator_grid_era5} and~\ref{fig:operator_grid_gs}
show what that looks like: ten panels, ten surfaces climbing monotonically to
the top row, not one of them turning over inside the grid. Extending the grid does not fix this: run to 2160~h the optima
slide to 720--1440~h, where a ``filtered driver'' is an annual cycle and its
correlation with the response is a statement about season. The negative control
does exactly the same, $r = -0.47$ at the bound, which is the tell---a channel
with no mechanism cannot acquire one by being smoothed, so what the smoothing
found is variation the whole site shares.

The consequence for the report is a restriction on what the levels are allowed
to say. On that band the study reports the instantaneous association and the two
identified operators, and treats a bound-pinned time constant as missing rather
than as measured. The diurnal band carries the coupling result.

\paragraph{This also corrects the source of the idea.} The retired study in
\texttt{studies/obsolete/inclination\_prediction/}, whose figures prompted this
scan, capped its own grid at 168~h and reported radiation gaining
$\Delta R^2 = 0.43$ from a time constant sitting exactly on that cap. That is
the same slide, unremarked. The method it demonstrates is sound and this report
adopts it; the number it produced is an artefact of where its grid ended.

\subsection{How to read one of these grids}
\label{sec:reading}

The three figures that follow are the study's only two-dimensional ones and
they carry most of what the tables cannot. This subsection says how to read
them.

\paragraph{What is drawn.} Each cell is one operator applied to the driver
before the driver is regressed on the response, and the two axes are the
operator's two parameters. They stand for the two distinct things a wall can do
to a forcing before the inclinometer sees it: answer late, or answer slowly.

The horizontal axis is the \emph{transport delay}, in hours. It shifts the
driver in time without changing its shape: at a delay of $d$ hours the response
now is paired with the driver as it was $d$ hours earlier. Physically it is
travel time. Heat entering the exposed face needs time to reach the depth whose
expansion governs the inclination, and the thermal front arrives there with its
shape intact but late. A delay is what a forcing looks like after crossing a
distance.

The vertical axis is the \emph{thermal time constant} $\tau$, in hours. It does
not shift the driver; it smooths and damps it. The driver is passed through a
first-order low-pass filter, the discrete form of a lumped thermal mass: each
hour the filtered value moves towards the current driver by the fraction
$\Delta t/(\tau+\Delta t)$ of the distance still between them, so that a step
in the driver is answered by an exponential approach that has covered about
63\,\% of its way after $\tau$ hours and 95\,\% after $3\tau$. A body of time
constant $\tau$ does not follow the forcing around it; it integrates it, with
a memory of about $\tau$ hours, and its temperature is a rounded copy of the
forcing whose fast swings are damped and whose peak arrives later the larger
$\tau$ is. A time constant of zero returns the driver untouched. The filter is
applied first and the delay second, following the physics: the mass integrates
the forcing, and the resulting thermal state then takes time to reach the depth
that moves the instrument.

The two are easy to confuse on a single sinusoid and separable on a real
record, which is why the grid scans them jointly rather than one at a time. On
a pure 24-hour cycle a delay of $d$ hours and a time constant whose phase lag
is $d$ hours are the same phase shift, and no correlation can tell them apart.
What separates them is that a delay moves every frequency by the same time,
whereas the filter lags and damps the fast components more than the slow ones:
on the diurnal cycle with its harmonics, the multi-day weather swings and the
sharp events such as a cold front or a clear day after cloud, a delayed driver
stays sharp where a filtered one is blurred, and the regression can tell which
shape the response follows. A scan over delay alone reads the two as one
number; the grid reads them as a pair.

The colour of a cell is the $R^2$ of the driver, put through that delay and
that time constant, regressed on the response over the hours the two series
share at that operator. Dark is a high $R^2$, following the project's rule that
the meaningful end of a scalar map is the dark one. The dot marks the winning
cell and the panel title repeats its coordinates and its $R^2$. The bottom-left
cell, $(0,0)$, is the instantaneous case: no delay, no inertia, the driver
exactly as recorded. Moving right along the bottom row asks whether the wall
answers late but sharply; moving up the left column asks whether it answers on
time but blurred; the interior asks both at once.

\paragraph{Two things the panels do not say, and one they do.} The vertical
axis is a list of chosen time constants drawn at equal spacing, not a linear
scale, so a step from 48 to 72~h occupies the same height as one from 0 to 1~h
and vertical distance is not proportional to hours. And $R^2$ discards the sign:
a panel cannot distinguish a driver that heats the wall from one that cools it,
and the sign must be read from the $r$ column of Table~\ref{tab:level}
or~\ref{tab:diurnal}.

What the panels do say, and this is deliberate, is how the drivers compare. All
three figures are drawn on one colour scale, running from 0 to the largest
$R^2$ any of their panels reaches, rounded up---0.90, set by the on-structure
air temperature. A shade therefore means the same number in every panel of
every one of the three figures, and nothing is clipped. The price is that a weak
panel occupies only the bottom of the range and shows little of its own
structure: the station's pressure is a uniform pale field because it explains
between $0.6$ and $1.0$~per cent of the variance everywhere on its grid, and
that flatness is the finding rather than a defect of the drawing.

\paragraph{The shapes, and what each one means.} A grid is read by its shape
before its maximum, because the shape says whether the maximum means anything.

\begin{itemize}
\item \textbf{A bright corner at $(0,0)$, fading in both directions.} The
coupling is instantaneous: neither a delay nor an inertia improves it. The
on-structure air temperature is this, and it is the cleanest result in the
study.

\item \textbf{A bright band along the bottom row, away from zero.} A pure
transport delay: the response reproduces the driver's shape, later. Nothing in
this study is cleanly of this kind, which is itself informative.

\item \textbf{A bright band up the left column, at zero delay.} A pure
inertia: the response follows a smoothed driver rather than a shifted one. The
on-structure radiation on the levels is this, rising from $R^2 = 0.166$ at
$\tau = 0$ to $0.288$ at $\tau = 48$~h. Note how weakly the maximum is placed
even so---$0.285$, $0.288$ and $0.287$ at 36, 48 and 72~h---so the shape is
established and the exact time constant is not.

\item \textbf{A diagonal ridge.} The two parameters are trading off: a longer
delay with a shorter time constant fits about as well as the reverse, because
both mostly move the driver's phase. Only the combination is resolved, and the
winning pair is one point on a ridge rather than a measurement of either
parameter. Wall temperature shows this around its optimum, where the four best
cells---$(-4, 4)$, $(-4, 3)$, $(-5, 4)$ and $(-5, 6)$---lie within $0.002$ of
each other in $R^2$.

\item \textbf{A surface that climbs to the top row and stops there.} The
optimum is the bound, not a turning point: the scan was stopped by the grid
rather than by the physics, and the time constant is unidentified. All ten
external panels of Figures~\ref{fig:operator_grid_era5}
and~\ref{fig:operator_grid_gs} are this shape. Such a cell must not be quoted
as a measured time constant, and the tables flag it.

\item \textbf{Two bright regions about twenty-four hours apart.} Aliasing. On
a near-periodic signal a delay of $d$ and a delay of $d \pm 24$ produce almost
the same alignment, so a scan wide enough to contain both will report whichever
happens to score higher. Wall temperature, scanned over signed delays, peaks at
$-4$~h with $R^2 = 0.802$ and again at $+19$~h with $0.777$: the second is the
first wearing a different sign of delay, and an unbounded scan that returned it
would describe the wall as following its own probe by most of a day.

\item \textbf{A pale, featureless panel.} No coupling at any operator, and on
the shared scale that pallor is directly comparable with every other panel
rather than relative to the panel's own contents. Station pressure is this, and
so, more surprisingly, is the on-structure radiation on the levels: its interior
maximum is real but it sits at $R^2 = 0.29$, against $0.89$ for the air
temperature beside it.
\end{itemize}

\paragraph{Reading a panel with the tables.} The figure locates the optimum;
two columns of the tables say whether it is worth anything.
$\Delta R^2$ is the winning cell's $R^2$ minus the instantaneous cell's, so a
small value means the operator is decoration however confident the panel title
looks---the on-structure air temperature's is exactly zero. For a driver scanned
over signed delays, $R^2$ lost to causality is what disappears when the search
is restricted to the non-negative half of the panel, which is the half a
forecast is allowed to use.

\paragraph{One panel, end to end.} Take the wall temperature in
Figure~\ref{fig:operator_grid}. The bright region is a ridge, so the delay and
the time constant are not separately sharp. Its maximum sits at a \emph{negative}
delay, so the response leads the probe, which Section~7.5 reads as a probe
mounted deeper than the layers that drive the movement. There is a second bright
region near $+19$~h, which is that maximum's daily alias and not a second
mechanism. The maximum is inside the grid on both axes, so unlike every external
driver this time constant is identified. Table~\ref{tab:level} then supplies
what the panel discarded: the association is negative, $r = -0.895$, the gain is
$-2.95$~mdeg per degree, the operator is worth $\Delta R^2 = 0.071$ over the
instantaneous fit, and restricting it to causal delays costs $0.024$.

\begin{figure}[H]
\centering
\includegraphics[width=\textwidth]{TR_F07_operator_grid_str.png}
\caption{The delay-by-time-constant grid for the four on-structure drivers with
a mechanism, on the levels, current era. Each panel is titled with its own
winning cell, marked in the accent colour. All three grid figures share one
colour scale, 0 to 0.90, taken from the strongest panel among them so that
nothing is clipped and a shade means the same number throughout. What the table
cannot show is here: where the bright region is a ridge rather than a
peak, the delay and the time constant are trading off against one another and
the winning pair is one point on that ridge. The wall temperature's second ridge,
around a delay of twenty hours, is the 24-hour alias of its optimum at $-4$ and
is what the signed scan's bound exists to keep out of the answer.}
\label{fig:operator_grid}
\end{figure}

The same grid for the two external source families says something the
on-structure one does not, and it says it at a glance.

\begin{figure}[H]
\centering
\includegraphics[width=\textwidth]{TR_F08_operator_grid_era5.png}
\caption{The grid for the ERA5 channels, on the levels, whole record. Every
panel has the same shape: the surface climbs monotonically with the time
constant and never turns over, so its maximum lies on the top row of the grid
rather than inside it. That is the visible form of a time constant the record
does not identify---the scan is not finding an optimum, it is being stopped by
the bound. Compare the on-structure panels of
Figure~\ref{fig:operator_grid}, where the wall temperature and the radiation
have an interior maximum with structure around it.}
\label{fig:operator_grid_era5}
\end{figure}

\begin{figure}[H]
\centering
\includegraphics[width=\textwidth]{TR_F09_operator_grid_gs.png}
\caption{The grid for the Gubbio station channels, on the same terms and the
same shared colour scale. The picture is the reanalysis's: five channels, five
maxima on the bound. Because the scale is shared with
Figure~\ref{fig:operator_grid}, the comparison between the families can be made
by eye---these panels are pale where the on-structure ones are dark, and the
station's pressure, coupled to nothing, is the palest field in the study.}
\label{fig:operator_grid_gs}
\end{figure}

% ===========================================================================
\section{Across the rounds}
\label{sec:across}

The point of screening three source families rather than one is that the three
disagree, measurably so in Study~02, and a coupling that survives the
disagreement is a statement about the wall rather than about a file.

Air temperature survives it. It is the leading or near-leading driver in every
round, always at zero lag on the diurnal band, always negative, with gains of
$-2.79$, $-2.23$ and $-2.04$~\mdegC{} for the housing, the station and ERA5
respectively. The ordering is the exposure ordering: the channel closest to the
wall gives the largest gain, and the spatial average the smallest.

Solar radiation survives it too, at zero to one hour on the diurnal band, with
gains between $-0.026$ and $-0.035$~mdeg per watt per square metre. On the
levels all three radiation channels want a long time constant and two of the
three want more than the grid allows, which Section~\ref{sec:operator} reads as
a seasonal covariation rather than a thermal one.

Relative humidity appears in every round with a \emph{positive} association of
$r \approx +0.69$ to $+0.79$. This is not an independent finding. Humidity at
this site is largely the inverse of the daily temperature cycle, and its
correlation with the response carries the sign that inversion implies. It is
reported because the screen reports what it measures, not because a
moisture mechanism has been demonstrated.

The tables below give the full screen, one row per driver, pooled over the whole
record.

\begin{table}[H]
\centering
\small
\begin{tabular}{lrrrrrrll}
\toprule
\textbf{Driver} & \textbf{Lag} & \textbf{$\tau$} & \textbf{r} &
\textbf{$\Delta R^2$} & \textbf{Slope} & \textbf{N} & \textbf{Clears} &
\textbf{Sign} \\
\midrule
\input{../outputs/TR_T03_coupling_level.tex}\\
\bottomrule
\end{tabular}
\caption{The screen on the levels. \textbf{Lag} and \textbf{$\tau$} in hours,
\textbf{$\Delta R^2$} what the operator bought over the instantaneous fit,
\textbf{Slope} in millidegrees per unit of the driver, \textbf{Clears} whether
the driver exceeds the negative control, \textbf{Sign} whether its slope points
the way the site's geometry predicts. A $\tau$ of 72~h is the grid's bound and
means the time constant was not identified; see Section~\ref{sec:operator}.}
\label{tab:level}
\end{table}

\begin{table}[H]
\centering
\small
\begin{tabular}{lrrrrrrll}
\toprule
\textbf{Driver} & \textbf{Lag} & \textbf{$\tau$} & \textbf{r} &
\textbf{$\Delta R^2$} & \textbf{Slope} & \textbf{N} & \textbf{Clears} &
\textbf{Sign} \\
\midrule
\input{../outputs/TR_T04_coupling_diurnal.tex}\\
\bottomrule
\end{tabular}
\caption{The same screen on the diurnal band, where the time constant is fixed
at zero because it is not identifiable against a single period. The band is
where the daily forcing lives, and where every association in this study is
strongest.}
\label{tab:diurnal}
\end{table}

\begin{figure}[H]
\centering
\includegraphics[width=\textwidth]{TR_F05_leading_driver_scatter.png}
\caption{The response against the strongest surviving driver at the lag the scan
selected, diurnal band, with the fitted gain and its interval.}
\label{fig:scatter}
\end{figure}

% ===========================================================================
\section{Stability}
\label{sec:stability}

A gain fitted over a whole record is a single number that a single season can
produce on its own. Re-fitting the same lag month by month separates a coupling
from a coincidence: a real one holds its slope, within its uncertainty, across
months in which the driver's own range changes.

The on-structure air temperature holds a median gain of $-2.48$~\mdegC{} across
thirty-five months, ranging between $-3.63$ and $-1.88$. The two external air
temperatures behave the same way about their own medians. Solar radiation is
stable to within a factor of about three across the same window. The
wall-temperature and on-structure radiation gains rest on nine months only,
which is the whole current era, and their stability is correspondingly weakly
established.

Wind speed channels are visibly unstable: ERA5 wind speed ranges from $-0.84$ to
$-14.6$~mdeg per m/s and station wind speed from $0.09$ to $3.02$ across months.
Neither should be carried forward as a driver on the strength of a pooled fit.

\begin{figure}[H]
\centering
\includegraphics[width=\textwidth]{TR_F06_gain_stability.png}
\caption{Gain re-fitted month by month at the lag the pooled scan selected, with
its Newey--West interval. One panel per driver: the gains are in different units
and an axis holding several of them would say nothing about any.}
\label{fig:stability}
\end{figure}

% ===========================================================================
\section{The control is not silent}
\label{sec:control}

The supply voltage was carried as a negative control on the reasoning that no
mechanism moves a wall by battery voltage, so whatever it scored would be the
study's noise floor. It did not behave as a null channel. On the levels over the
post-outage record it reaches $r = -0.472$, and on the diurnal band it reaches
$-0.673$ within the post-outage legacy era and $+0.549$ within the current one.

On the levels it also takes the largest time constant the grid offers, landing
at $r = -0.47$ against the bound exactly as the candidate drivers do.

The explanation is not that the wall responds to the battery. It is that the
battery responds to the same daily cycle as everything else---a solar-charged
supply in a sun-exposed housing carries the day in its voltage---and, on the
slow band, to the same annual one. The control sliding to the bound alongside
the candidates is what identifies that bound as an artefact rather than a
measurement, and it is the clearest thing the control did in this study. The control is
therefore not mechanism-free, and the floor it sets is a \emph{conservative}
one---too high rather than too low. A driver that clears it has cleared a
channel that itself carries the diurnal forcing, which makes clearing a stronger
statement than it was designed to be; a driver that fails to clear it has not
thereby been shown to be uncoupled.

Read that way, thirteen of the twenty-two screened drivers clear the control on
the diurnal band; the channels that do not include precipitation, the wind
direction components, station pressure, and ERA5 dew point.

A better control for a future study would be a channel with a comparable
sampling history and no thermal exposure at all. The project does not currently
record one.

% ===========================================================================
\section{Verdict}
\label{sec:verdict}

\begin{table}[H]
\centering
\small
\begin{tabular}{llrrrrrrlll}
\toprule
\textbf{Band} & \textbf{Driver} & \textbf{Lag} & \textbf{$\tau$} &
\textbf{r} & \textbf{$\Delta R^2$} & \textbf{Slope} & \textbf{N} &
\textbf{Clears} & \textbf{Signif.} & \textbf{Sign} \\
\midrule
\input{../outputs/TR_T06_verdict.tex}\\
\bottomrule
\end{tabular}
\caption{The verdict, pooled over the whole record. The three conditions are
reported separately rather than combined, because a driver can fail one and pass
another and the distinction matters to what a later study may do with it.}
\label{tab:verdict}
\end{table}

What a later study may take from this report:

\begin{itemize}
\item \textbf{Air temperature is the driver.} Negative, at zero lag on the
diurnal band, in all three source families, with a gain between $-2.0$ and
$-2.8$~\mdegC{} that is stable over thirty-five months. Where the
on-structure channel is present it is the strongest of the three; where it is
absent, either external source carries the same sign and lag at a reduced gain.

\item \textbf{Solar radiation couples at zero to one hour} with a gain of order
$-3\times10^{-2}$~mdeg per watt per square metre, and the external radiation
channels reproduce the on-structure one closely enough to stand in for it
outside the certified window.

\item \textbf{Wall temperature leads nothing.} The inclination leads it by two
hours on the diurnal band and by four, with a four-hour time constant, on the
levels. It is the only driver whose time constant this record identifies. Use
the measurement as a diagnostic of probe depth; clamping the lead to zero, which
any forecast requires, costs $\Delta R^2 = 0.068$ on the band and $0.024$ on the
levels.

\item \textbf{No thermal inertia is identified for the external drivers.} Their
level-band time constants pin against the grid's bound and are reported as
missing. A model built on this study should carry them instantaneously, at the
delays the diurnal band measured, and should not inherit a time constant from
Table~\ref{tab:level}.

\item \textbf{The sign agrees with the site's structural convention} for every
heating driver in every round, on the compensated response. This does not
resolve the contradiction of Section~7.5, which concerns the raw channel, and it
must not be quoted as though it did.
\end{itemize}

% ===========================================================================
\section{Limitations}
\label{sec:limitations}

\begin{itemize}
\item \textbf{The compensated response is not an independent measurement of the
structure.} Section~7.2's compensation subtracts a fixed
5~\mdegC{} from a channel whose own thermal slope is of comparable size, so
part of what is measured here as a structural response is the residual of that
subtraction. Screening the raw channel alongside it was deliberately left out of
scope; it is the obvious next study.

\item \textbf{Zero lag at hourly resolution is a bound, not a value.} A response
that follows its driver by twenty minutes and one that follows it instantly are
the same number on this grid. Study~01's archive is on a twenty-minute grid and
could support a finer scan.

\item \textbf{The levels cannot separate inertia from season.} Section
\ref{sec:operator} sets this out: fifteen of twenty-two time constants pin
against the grid's bound, and the negative control pins with them. A record of
three years contains only three realisations of the annual cycle, which is not enough
to tell a slow thermal response from a seasonal coincidence. Nothing short of a
different experiment---or a driver whose annual cycle differs from the
response's---will settle it.

\item \textbf{The current-era channels rest on nine months.} Wall temperature
and on-structure radiation carry roughly five thousand paired hours against
twenty-one thousand for air temperature, and their stability is established
correspondingly weakly.

\item \textbf{The negative control is thermally exposed}, as
Section~\ref{sec:control} sets out. The floor it defines is conservative and the
screen inherits that.

\item \textbf{Correlation between drivers is not addressed.} Air temperature,
dew point, humidity and radiation share one daily cycle, and this study screens
each against the response separately rather than fitting them together. Which of
them a model should carry, and with what collinearity treatment, is a question
for the modelling study, not for this screen.
\end{itemize}

\end{document}
